\section{Closed-form searches and quadratic conjugacy}
\label{a:app:closed-form-search}

This appendix collects the negative and structural calculations that help interpret
\(A(p)\) but are not needed for its probabilistic role in the main text.  The
numerical formulae are approximations only.  Their value is diagnostic: they show
how little a visually accurate fit to \(A(p)\) determines the singular quantity
\(1/A(p)\).

\subsection{Exploratory formulae}

Polynomial and rational searches can be constrained to the known endpoints
\(A(0)=1/2\) and \(A(1/2)=0\).  One of the more economical rational candidates was
\begin{equation}
  \widehat A_1(p)
  =
  \frac{\tfrac12(1-2p)(1+2p)}
       {1+\tfrac12p-\tfrac52p^2-\tfrac65p^3+\tfrac65p^4}.
  \label{a:eq:Ahat1}
\end{equation}
The candidate passes through both endpoints.  At the critical endpoint, however,
\begin{equation}
  \widehat A_1'(1/2)=-\frac{2}{0.55}\approx-3.636,
\end{equation}
whereas \cref{a:thm:A-near-critical} requires the left derivative \(-4\).
Matching the endpoint values while missing this slope produces an unacceptable
relative error in \(A\), and hence in \(1/A\), near criticality.

A less restricted symbolic tree search produced many longer formulae with similar
visual accuracy.  Two representative outputs were
\begin{multline}
  \widehat A_2(p)
  =0.31696448624690005
  \log\!\Biggl(
  \Biggl|\cos\!\biggl(
  \left(\frac{1}{3.19801133824532}\right)^{1/4}\\
  {}+p\biggl[
  \sin\!\Bigl(
    \cos\!\bigl(|p|(0.621865696665606-p)
    \cos(\e^{\e^p})\bigr)
  \Bigr)+|p|
  \biggr]\biggr)\Biggr|
  \Biggr)
  +0.5982282661889831
  \label{a:eq:Ahat2}
\end{multline}
and
\begin{equation}
  \widehat A_3(p)
  =
  \frac{-0.616647881739505\,
  \e^{-0.552605335919693\,\e^p}}
  {\sin(\cos p)-|p|}
  +0.921964323679771.
  \label{a:eq:Ahat3}
\end{equation}
These expressions have no inferential status.  They record a curve-fitting
exercise and are not proposed identities.

\begin{figure}[htbp]
  \centering
  \includegraphics[width=0.60\textwidth]{A3_hat_plot.pdf}
  \caption{The fitted curve \(\widehat A_3\) and product evaluations of \(A(p)\).
  The discrepancy is barely visible on this scale, although
  \(\widehat A_3(1/2)\approx9.1\times10^{-4}\) rather than zero and its endpoint
  slope is approximately \(-3.63\) rather than \(-4\).}
  \label{a:fig:Ahat3}
\end{figure}

One can force the correct critical branch by splicing:
\begin{equation}
  \widehat A_4(p)
  =
  \begin{cases}
  -0.04011483983619545\,\e^{\e^{2p}}
    +0.6079994336528085,
    &p<0.499962,\\[1mm]
  2-4p,
    &p\geq0.499962.
  \end{cases}
  \label{a:eq:Ahat4}
\end{equation}
This construction is a practical surrogate rather than a closed form for \(A\).
If a splice is needed, the convergent product in \cref{a:eq:prodFormula} is
preferable on the first branch and the proved asymptotic
\cref{a:eq:A-near-critical} on the second.

\subsection{Conjugacy with the quadratic family}
\label{a:app:quadratic-conjugacy}

The logistic family \(f_r(x)=rx(1-x)\) is affinely conjugate to the quadratic
family
\begin{equation}
  P_c(z)=z^2+c.
\end{equation}
Set
\begin{equation}
  z=r\left(\frac12-x\right).
  \label{a:eq:logistic-affine-coordinate}
\end{equation}
If \(x'=f_r(x)\) and \(z'=r(1/2-x')\), direct substitution gives
\begin{equation}
  z'=z^2+c,
  \qquad
  c=c(r):=\frac{2r-r^2}{4}.
  \label{a:eq:logistic-quadratic-conjugacy}
\end{equation}
Thus the subcritical branching window \(0<r<1\) maps onto \(0<c<1/4\).

\begin{figure}[htbp]
  \centering
  \includegraphics[width=0.68\textwidth]{figure1_parameter_conjugacy.pdf}
  \caption{The parameter correspondence
  \(c(r)=(2r-r^2)/4\).  The subcritical window \(0<r<1\) maps to
  \(0<c<1/4\); the elementary monomial and Chebyshev parameters described below
  lie outside that window.}
  \label{a:fig:parameter-conjugacy}
\end{figure}

The interval \(0<c<1/4\) is the real radius of the main cardioid of the
Mandelbrot set.  Indeed, a fixed point of \(z^2+c\) with multiplier \(r\) has
coordinate \(z_\ast=r/2\), and its fixed-point equation gives
\begin{equation}
  c=z_\ast-z_\ast^2
   =\frac r2-\frac{r^2}{4}.
\end{equation}
Thus the endpoint \(r\uparrow1\) approaches the parabolic cusp \(c=1/4\).
This observation locates the parameter range but supplies no evidence for a
closed-form claim.

\begin{figure}[htbp]
  \centering
  \includegraphics[width=0.88\textwidth]{figure4_mandelbrot_context.png}
  \caption{The image \(0<c<1/4\) of the subcritical branching interval within
  the quadratic parameter plane.  The exceptional parameters \(c=0\) and
  \(c=-2\) are shown for comparison.}
  \label{a:fig:mandelbrot-context}
\end{figure}

\subsection{The two elementary comparison cases}
\label{a:app:elementary-linearizers}

The exceptional quadratic maps conjugate to a monomial or Chebyshev polynomial
have elementary linearisers.  They provide useful algebraic checks, although
neither lies in \(0<r<1\).

For \(r=2\), put \(u=1-2x\).  Then
\begin{equation}
  1-2f_2(x)=(1-2x)^2=u^2,
\end{equation}
and the normalised local lineariser is
\begin{equation}
  \psi_2(x)=-\frac12\log(1-2x).
  \label{a:eq:psi-r2}
\end{equation}
It satisfies
\(\psi_2(f_2(x))=2\psi_2(x)\),
\(\psi_2(0)=0\) and \(\psi_2'(0)=1\).

For \(r=4\), write \(x=\sin^2\theta\).  The double-angle identity gives
\begin{equation}
  f_4(x)=4x(1-x)=\sin^2(2\theta),
\end{equation}
so that
\begin{equation}
  \psi_4(x)=\{\arcsin(\sqrt{x})\}^2
  \label{a:eq:psi-r4}
\end{equation}
satisfies
\(\psi_4(f_4(x))=4\psi_4(x)\) and
\(\psi_4(x)\sim x\) at the origin.

The distinction is categorical.  At \(r=2\) and \(r=4\), the origin has multiplier
greater than one, and the formulae \cref{a:eq:phi2,a:eq:phi4} are local repelling
linearisers.  In the branching window, the origin is attracting and
\cref{a:eq:schroder} defines the attracting Koenigs coordinate.  The elementary
formulae are therefore comparison cases, not analytic continuations of \(A(p)\).

Finally, \cref{a:thm:main-ht} concerns \(z\mapsto\psi_r(z)\) at fixed \(r\); it
does not prove that the parameter evaluation \(p\mapsto2\psi_{2p}(1/2)\) is
non-elementary.  Maintaining this separation is the substantive correction to the
stronger claim suggested by the original closed-form search.

\FloatBarrier
